\kchapter{Pods: The Smallest Unit}{pod}

\begin{objectives}
\begin{itemize}
  \item what a Pod is and why Kubernetes does not run bare containers;
  \item the Pod lifecycle, phases, and restart policies;
  \item commands, arguments, and environment variables;
  \item multi-container Pods: init containers and sidecars;
  \item health checks: liveness, readiness, and startup probes;
  \item graceful shutdown.
\end{itemize}
\end{objectives}

\section{What is a Pod?}

A \kterm{Pod} is a group of one or more containers that are always scheduled
together on the same node and share:

\begin{itemize}
  \item one \textbf{network identity} --- a single IP address and port space
        (containers in a Pod reach each other on \kc{localhost});
  \item \textbf{storage volumes} that any container in the Pod can mount;
  \item a \textbf{lifecycle} --- they are created and destroyed together.
\end{itemize}

\begin{center}
\begin{tikzpicture}[
  ctr/.style={fill=k8sblue, text=white, rounded corners=4pt, minimum width=2.6cm,
              minimum height=1.1cm, font=\small\bfseries, align=center}]
  \node[ctr] (a) at (0,0) {\faCube\ \ app\\[-2pt]\scriptsize\mdseries :8080};
  \node[ctr, fill=k8sdeep] (b) at (3.4,0) {\faCube\ \ log-shipper\\[-2pt]\scriptsize\mdseries reads /logs};
  \node[draw=k8smid, fill=k8slight, rounded corners=3pt, minimum width=6cm,
        font=\scriptsize\bfseries, text=k8snavy] (v) at (1.7,-1.4) {\faHdd\ \ shared volume: /logs};
  \begin{scope}[on background layer]
    \node[draw=k8sblue, line width=1.3pt, fill=white, rounded corners=10pt,
          inner sep=0.4cm, fit=(a)(b)(v)] (pod) {};
  \end{scope}
  \node[anchor=south west, font=\small\bfseries, text=k8sblue] at (pod.north west)
       {\kicon{pod}\ Pod \quad\mdseries\footnotesize IP 10.244.1.17};
  \draw[<->, thick, k8sorange] (a) -- node[above, font=\scriptsize, text=k8sorange]{localhost} (b);
\end{tikzpicture}
\end{center}

\begin{analogy}
A Pod is like a \emph{shared office}. The people (containers) inside share one
street address (IP), one mailbox (port space), and filing cabinets (volumes).
Most offices have a single occupant; occasionally an assistant moves in.
\end{analogy}

\begin{keyidea}
Pods are \textbf{ephemeral and disposable}. They are not repaired; they are
replaced. A replaced Pod gets a new name and a new IP. That is why you will
almost never create Pods directly --- you will use a Deployment (Chapter 8)
and a Service (Chapter 10).
\end{keyidea}

\section{Your first real Pod}

\begin{yamlcode}{hello-pod.yaml}
apiVersion: v1
kind: Pod
metadata:
  name: hello
  labels:
    app: hello
spec:
  containers:
    - name: web
      image: nginx:1.27
      ports:
        - containerPort: 80
          name: http
      env:
        - name: GREETING
          value: "Hello, Kubernetes"
      resources:              # covered in depth in Chapter 14
        requests:
          cpu: 50m
          memory: 64Mi
        limits:
          memory: 128Mi
\end{yamlcode}

\begin{bashcode}{terminal}
kubectl apply -f hello-pod.yaml
kubectl get pod hello -o wide
kubectl describe pod hello      # read the Events at the bottom
kubectl exec -it hello -- sh -c 'echo $GREETING'
kubectl delete pod hello
\end{bashcode}

\section{Command, args, and environment}

Containers start with the image's \kc{ENTRYPOINT} and \kc{CMD}. In Kubernetes
you can override them: \kc{command} replaces \kc{ENTRYPOINT}, \kc{args}
replaces \kc{CMD}.

\begin{yamlcode}{busybox-pod.yaml}
apiVersion: v1
kind: Pod
metadata:
  name: counter
spec:
  restartPolicy: OnFailure
  containers:
    - name: counter
      image: busybox:1.36
      command: ["sh", "-c"]
      args: ["for i in 1 2 3 4 5; do echo tick $i from $POD_NAME; sleep 1; done"]
      env:
        - name: POD_NAME              # the Downward API: facts about the Pod itself
          valueFrom:
            fieldRef:
              fieldPath: metadata.name
        - name: NODE_NAME
          valueFrom:
            fieldRef:
              fieldPath: spec.nodeName
\end{yamlcode}

\section{The Pod lifecycle}

\begin{center}
\resizebox{\linewidth}{!}{%
\begin{tikzpicture}[
  ph/.style={rounded corners=12pt, minimum width=2.6cm, minimum height=1cm,
             font=\small\bfseries, text=white, align=center},
  >={Stealth[length=3mm]}]
  \node[ph, fill=k8sgray] (p) at (0,0) {Pending};
  \node[ph, fill=k8sblue] (r) at (4,0) {Running};
  \node[ph, fill=k8sgreen] (s) at (8,1.1) {Succeeded};
  \node[ph, fill=k8sred] (f) at (8,-1.1) {Failed};
  \node[ph, fill=k8sorange] (u) at (4,-2.2) {Unknown};
  \draw[->, thick] (p) -- node[above, font=\scriptsize]{scheduled +} node[below, font=\scriptsize]{images pulled} (r);
  \draw[->, thick] (r) -- node[above, sloped, font=\scriptsize]{all exit 0} (s);
  \draw[->, thick] (r) -- node[below, sloped, font=\scriptsize]{a container fails} (f);
  \draw[->, thick, dashed] (r) -- node[right, font=\scriptsize]{node lost} (u);
\end{tikzpicture}}
\end{center}

\begin{center}
\ktablestart
\begin{tabularx}{\linewidth}{P{2.4cm} L}
\kheader \kth{Phase} & \kth{Meaning} \\
Pending & Accepted, but not yet running: waiting to be scheduled or pulling images. \\
Running & Bound to a node; at least one container is running. \\
Succeeded & All containers exited successfully and will not restart. \\
Failed & All containers have stopped and at least one failed. \\
Unknown & The node stopped reporting. \\
\end{tabularx}
\end{center}

The \kc{STATUS} column of \kc{kubectl get pods} often shows more detailed
reasons. Learn these by heart:

\begin{center}
\ktablestart
\begin{tabularx}{\linewidth}{P{3.6cm} L}
\kheader \kth{Status} & \kth{Usual cause} \\
\kc{ContainerCreating} & Pulling image, mounting volumes. Normal for a few seconds. \\
\kc{ImagePullBackOff} & Wrong image name/tag, private registry without credentials. \\
\kc{CrashLoopBackOff} & The container keeps exiting. Check \kc{kubectl logs --previous}. \\
\kc{OOMKilled} & Exceeded its memory limit. \\
\kc{Pending} (long) & No node has enough CPU/memory, or a PVC is unbound. Check events. \\
\kc{CreateContainerConfigError} & Referenced ConfigMap/Secret/key does not exist. \\
\end{tabularx}
\end{center}

\subsection{Restart policy}

\kc{spec.restartPolicy} tells the kubelet what to do when a container exits:
\kc{Always} (default; for long-running services), \kc{OnFailure} (for batch
jobs), or \kc{Never}. Restarts use exponential back-off: 10\,s, 20\,s, 40\,s
\ldots up to 5 minutes --- hence ``CrashLoop\emph{BackOff}''.

\section{Multi-container Pods}

\subsection{Init containers}

\kterm{Init containers} run \emph{to completion, in order}, before the main
containers start. Use them to wait for a dependency, run migrations, or
prepare files.

\begin{yamlcode}{init-demo.yaml}
apiVersion: v1
kind: Pod
metadata:
  name: init-demo
spec:
  initContainers:
    - name: wait-for-db
      image: busybox:1.36
      command: ["sh", "-c", "until nc -z postgres 5432; do echo waiting; sleep 2; done"]
  containers:
    - name: app
      image: nginx:1.27
\end{yamlcode}

\subsection{Sidecar containers}

A \kterm{sidecar} runs alongside the main container for its whole life: a log
shipper, a proxy, a config reloader. Since Kubernetes 1.29 the recommended way
is a \emph{native sidecar}: an init container with \kc{restartPolicy: Always}.
It starts before the app, keeps running, and is stopped after the app.

\begin{yamlcode}{sidecar-demo.yaml}
apiVersion: v1
kind: Pod
metadata:
  name: sidecar-demo
spec:
  volumes:
    - name: logs
      emptyDir: {}
  initContainers:
    - name: log-shipper              # native sidecar
      image: busybox:1.36
      restartPolicy: Always
      command: ["sh", "-c", "touch /logs/app.log; tail -F /logs/app.log"]
      volumeMounts:
        - name: logs
          mountPath: /logs
  containers:
    - name: app
      image: busybox:1.36
      command: ["sh", "-c", "while true; do date >> /logs/app.log; sleep 2; done"]
      volumeMounts:
        - name: logs
          mountPath: /logs
\end{yamlcode}

\begin{bashcode}{terminal}
kubectl apply -f sidecar-demo.yaml
kubectl logs sidecar-demo -c log-shipper -f
\end{bashcode}

\section{Health checks: probes}

The kubelet can check your container's health in three different ways, each
with a different consequence:

\begin{center}
\ktablestart
\begin{tabularx}{\linewidth}{P{2.4cm} L L}
\kheader \kth{Probe} & \kth{Question it asks} & \kth{If it fails\ldots} \\
\textbf{startup} & ``Has the app finished starting?'' &
   Other probes wait. After \kc{failureThreshold}, the container is restarted. \\
\textbf{liveness} & ``Is the app still alive, or stuck?'' &
   The container is \textbf{restarted}. \\
\textbf{readiness} & ``Can the app serve traffic right now?'' &
   The Pod is \textbf{removed from Service endpoints} (no restart). \\
\end{tabularx}
\end{center}

Probes can use \kc{httpGet}, \kc{tcpSocket}, \kc{grpc}, or \kc{exec} (run a
command; exit code 0 means healthy).

\begin{yamlcode}{probes.yaml}
containers:
  - name: api
    image: ghcr.io/example/api:1.4.0
    ports:
      - containerPort: 8000
    startupProbe:
      httpGet: { path: /healthz, port: 8000 }
      periodSeconds: 5
      failureThreshold: 24       # allow up to 2 minutes to start
    livenessProbe:
      httpGet: { path: /healthz, port: 8000 }
      periodSeconds: 10
      failureThreshold: 3
    readinessProbe:
      httpGet: { path: /readyz, port: 8000 }
      periodSeconds: 5
      failureThreshold: 2
\end{yamlcode}

\begin{warningbox}
Do \textbf{not} check your database inside the \emph{liveness} probe. If the
database blips, every replica fails liveness and Kubernetes restarts your whole
fleet --- turning a small outage into a big one. Dependencies belong in
\emph{readiness}; liveness should only answer ``is this process wedged?''.
\end{warningbox}

\section{Graceful shutdown}

When a Pod is deleted (scale-down, rollout, node drain), Kubernetes:

\begin{enumerate}
  \item marks the Pod \kc{Terminating} and removes it from Service endpoints;
  \item runs any \kc{preStop} hook;
  \item sends \kc{SIGTERM} to each container's main process;
  \item waits up to \kc{terminationGracePeriodSeconds} (default 30\,s);
  \item sends \kc{SIGKILL} to anything still running.
\end{enumerate}

\begin{yamlcode}{graceful.yaml}
spec:
  terminationGracePeriodSeconds: 45
  containers:
    - name: api
      image: ghcr.io/example/api:1.4.0
      lifecycle:
        preStop:
          exec:
            command: ["sh", "-c", "sleep 5"]   # let load balancers notice first
\end{yamlcode}

\begin{protip}
Make your app handle \kc{SIGTERM}: stop accepting new requests, finish
in-flight ones, close DB connections, exit. A short \kc{preStop} sleep avoids
dropped requests while endpoint removal propagates across the cluster.
\end{protip}

\section{Static Pods and Pod templates}

Two more terms you will see:

\begin{itemize}
  \item \kterm{Static Pods} are defined by files on a node (e.g.\
        \kc{/etc/kubernetes/manifests}) and run directly by the kubelet.
        That is how \kc{kubeadm} runs the control plane itself.
  \item A \kterm{Pod template} is the \kc{template:} section inside
        Deployments, Jobs, StatefulSets, and so on --- a blueprint those
        controllers stamp out Pods from.
\end{itemize}

\begin{recap}
\begin{itemize}
  \item A Pod = one or more containers sharing network, storage, and lifecycle.
  \item Pods are cattle, not pets: replaced, never repaired.
  \item Init containers prepare; native sidecars accompany; the app does the work.
  \item Startup, liveness, and readiness probes answer different questions.
  \item Handle \kc{SIGTERM} for zero-downtime shutdowns.
\end{itemize}
\end{recap}
